\clearpage
\chap{Extension proofs and their verification scope}\label{app:written}
The WR catalogue is deliberately separate from the eighty KS exposition items. It began with twelve complete written theorem groups. Version 1.6.0 added WR13--WR16 and supplied a focused formal counterpart for selected parts of the catalogue. The present edition expands all sixteen printed proofs while retaining those evidence boundaries. It does not promote a written limit theorem merely because a related pointwise estimate has now been compiled.

\begin{center}
\begin{tabularx}{\linewidth}{@{}p{.19\linewidth}X@{}}
\toprule Coverage & Exact interpretation\\\midrule
Complete focused & WR06, WR07, WR13 and WR14. The principal universal statements have typed counterparts in the three new modules and their actual protected bridge. Explanatory examples retain their stated separate scope.\\[5pt]
Partial focused & WR09 and WR16. The fixed-parameter and pointwise solver/data inequalities are formal. Their parameter-limit consequences still use written WR08.\\[5pt]
Written only & WR01--WR05, WR08, WR10--WR12 and WR15. Full proofs are printed. No new typed counterpart is assigned.\\\bottomrule
\end{tabularx}
\end{center}

The retained version 1.6.0 focused replay recompiled seven targets from source: three new modules, three inherited dependencies, and the root. It uses the pinned official Lean compiler and imported dependency caches. The canonical \path{formal/} tree remains unchanged, with its original accepted cumulative evidence. The source resolver binds WR references to the focused imported declaration types. This is an editorial correspondence, not automatic natural-language proof verification.

\subsection{Existence and budget geometry}
\textbf{WR01. Bounded approximate feasibility.} The real Hilbert operator and residual datum are fixed. A uniform correction-norm bound, not residual decay alone, gives an exact preimage by weak subsequence compactness. This remains a written proof.

\textbf{WR02. Nested affine minimum corrections.} Nonempty closed affine constraint sets are nested in one Hilbert space. The exact energy identity makes bounded minimum-norm corrections converge strongly. No formal compactness or affine-limit module is added here.

\textbf{WR03. Attained budget profile.} A finite nonnegative budget has an attained residual minimum even when the full operator range is not closed. The proof retains the weak topology argument rather than substituting finite-dimensional compactness.

\textbf{WR04. Sharp budget support.} The general support formula uses the actual restricted adjoint. The focused module proves the support inequality and supplies equality at a regularized optimizer's own budget. It does not prove the entire general budget-attainment route, so WR04 remains written-only as a group.

\textbf{WR05. Robust budget comparison.} The operators act on the same correction space, with operator and data errors bounded by $\eta$ and $\delta$. The bound $\delta+\eta M$ does not apply automatically when the protected kernel changes.

\subsection{Regularization and its limit}
\textbf{WR06. Constructed regularized minimizer.} Positive $\alpha$ makes $T^*T+\alpha I$ coercive. The formal proof constructs an actual continuous linear equivalence, defines the correction using its inverse, and proves the energy identity and uniqueness. The manuscript's self-adjoint closed-range proof and the formal coercive-form proof establish the same stated result. The bridge uses the existing \path{ProtectedHilbertRepair.restricted} operator.

\textbf{WR07. Parameter and budget comparison.} The formal statements compare two positive parameters, prove the two norm monotonicities, and show that no correction of smaller or equal norm has a lower residual than the regularized answer. The normalized residual has unit norm and attains the support expression whenever it is nonzero. The Lean formulation gives the optimizer and universal competitor inequality directly rather than introducing a second numerical profile definition.

\textbf{WR08. The regularization limit.} Projection onto the closure of the range is separated from actual range membership. Strong convergence of corrections requires the projected datum to be attained. Otherwise the correction norms diverge. These infinite-limit arguments remain written-only.

\textbf{WR09. Data sensitivity and parameter choice.} The focused proof supplies the coefficients $1/(2\sqrt\alpha)$ for correction noise and one for fitted-data noise. The scalar sharpness calculations and the parameter-choice convergence consequence retain their printed proofs. Formalizing a fixed-parameter inequality does not formalize its convergence application automatically.

\textbf{WR10--WR12. The explicit diagonal models.} The exact regularized family, omitted-tail bounds, matched-budget limit $\pi/4$, and attainable noisy datum remain written with exact finite enclosure tests. No general infinite-sum formalization or new diagonal Lean instance is claimed.

\subsection{Inexact solution and coordinate interpretation}
\textbf{WR13. A posteriori solver bounds.} The correction error, fitted error, and objective excess are bounded in terms of the actual normal residual and positive parameter. The focused theorem quantifies over every proposal in the Hilbert source. The protected bridge transfers its candidate-distance bound to the actual kernel subtype. Scalar sharpness examples are separated from the principal theorem and checked independently.

\textbf{WR14. A regularized primal and dual certificate.} An exact two-square identity proves weak duality, a dual maximizer constructed from the primal correction, global objective bracketing, and a correction-distance bound. The witness is arbitrary. Setting it to the proposal's fit residual gives the exact normal-residual gap.

\textbf{WR15. The coordinate covector metric.} A full independent kernel basis gives $g=BG^{-1}q$ and $\norm g^2=q^{\mathsf T}G^{-1}q$. The invariance proof and its invertibility hypotheses are printed and rationally tested. They have not been compiled as a new matrix theorem.

\textbf{WR16. A justified stopping rule.} The focused theorem proves the pointwise sum of the solver and data bounds. The printed convergence consequence uses written WR08 and the additional assumptions $\alpha\to0$, $\delta/\sqrt\alpha\to0$, and $\epsilon/\alpha\to0$. The displayed diagonal schedule also uses written WR12.

\subsection{Independent exact implementation and restart}
The rerun regularization suite contributes \EditionRegularizationChecks\ assertions and the rerun inexact-certificate suite \EditionInexactChecks. The independent oracle solves ambient constrained systems by determinant formulas rather than using the certificate program's kernel coordinates. The verifier checks a full kernel basis, retains its Gram metric, preserves the observable, and distinguishes measured fit from the actual residual. All outputs are exact rationals or squared bounds. It recomputes records and rejects altered claims without calling an exact optimizer.

